\subsection{第五係数非零の行列式零境界の全排除}
\label{subsec:cubic-CY-four-point-determinant}
\begin{theorem}\label{thm:cubic-CY-four-point-determinant-empty}
四単純基点の元表示で
\[
 x_0x_1x_5\ne0,\quad x_3=x_4=0,\quad x_0x_2-x_1^2=0
\]
を満たす全領域と、その全monomial軌道はH3包絡を与えない。
\end{theorem}
\begin{proof}
独立な対角重みで $x_0=x_1=x_5=1$ とすれば $x_2=1$ となる。
固定出典\cite{AS4CYpinned}の一般行列式境界の必要被覆FG1--FG5を用いる。
全21混合係数を残した元Schur行列を、$\lambda,\lambda+1$ だけを単元として消去する。
残余二式の和は
\[
 \frac{\lambda+1}{8}\{16(x_{15}-x_{16})^2+1\}
\]
である。従って $x_{15}-x_{16}=\pm i/4$。
まず正符号を扱い、$x_{16}=x_{15}-i/4$ として再び全単元消去を行う。
残る変数は $x_{12},x_{14},x_9$ であり、全残余の一式が単元倍を除き
\[
 x_9(16x_9-\lambda^2+1)=0
\]
となる。
当repoではこの二段階の全消去手順について、各式が実際の元残余式であること、
pivotが明記した単元であること、全残余式の一致を独立に確認した。

$x_9=(\lambda^2-1)/16$ の枝の全idealは、
$\lambda(\lambda+1)$ の局所化で
\[
 x_{12}=0,\qquad16(\lambda+1)^2x_{14}=\lambda^4
\]
と一致する。このideal一致を有理数体のガウス拡大上で再計算した。
全消去係数を戻すと、既存 $R^+_\lambda$ の元六関係と文字どおり一致する。
形式変数 $\lambda$ を固定して $i$ を反転することにより、負符号の全枝は
$R^-_\lambda$ に移る。いずれも命題\ref{prop:cubic-boundary-R}で全許容値が排除される。

$x_9=0$ の枝の全idealは、同じ局所化で
\[
 (\lambda-1)(\lambda^2+1)=0,\quad x_{12}(\lambda^2+1)=0,\quad
 16(\lambda+1)^2(x_{14}-x_{12}^2)=\lambda^2
\]
と一致する。従って $\lambda=1,i,-i$ だけを扱えばよい。
$\lambda=1$ では $x_{12}=0,x_{14}=1/64$ となり、全消去係数は
両符号それぞれ既存の $R^\pm_1$ に一致する。

$\lambda=\pm i$ では $x_{12}=h$ を任意に残し、
$x_{14}=h^2+\lambda/32$ となる。
$\lambda=i$ の全係数は既存 $J^\eta_h$ に一致する。
$\lambda=-i$ には
\[
 P=\operatorname{diag}(1,i,i,i)P_{(bc)}
\]
を作用させると、$\lambda=i$ の表示へ移る。
当repoでは両符号・両純虚値の四場合全てについて、元16語上の関係基底を可逆に移し、
$h$ を $\pm h,\pm ih$ のいずれかに替えた既存 $J$ 表示と一致することを確認した。
この移送に $h$ の除算はなく、$h=0$ も含まれる。
命題\ref{prop:cubic-boundary-J}は $h\ne0$ の全中心三次消滅と、
$h=0$ の五次商次元58による非AS4性の両方を既に証明している。
その既存の全値排除を適用すれば、この全境界もH3でない。
\end{proof}

公開検証器は\\
\texttt{verify\_cubic\_CY\_four\_point\_determinant.py}。
固定上流の一般境界と純虚境界の全control再実行、当repoの必要消去と元表示の独立移送、
既存 $R/J$ の全値排除を区別する。
AS十分性、全自己同型群、相対scheme分類は同時に輸入しない。
$x_5=0$ の行列式零境界、および出次数2以上の必要係数被覆は未完成である。
